% Part: normal-modal-logic
% Chapter: axioms-systems
% Section: consistency

\documentclass[../../../include/open-logic-section]{subfiles}

\begin{document}

\olfileid[fr]{nml}{prf}{con}

\olsection{Cohérence}

La cohérence est une propriété importante des ensembles de
!!{formula}s. Un ensemble de !!{formula}s est incohérent si l'on peut
en !!{derive} une contradiction, telle que~$\lfalse$ ; sinon, il est
cohérent. Si un ensemble est incohérent, ses !!{formula}s ne peuvent
pas être toutes vraies en un même monde d'un modèle où toutes les instances des axiomes du système sont vraies. Pour le théorème
de complétude, nous démontrerons la réciproque : tout ensemble cohérent
est vrai en un monde d'un modèle, à savoir le « modèle canonique ».

\begin{defn}
  Un ensemble $\Gamma$ est \emph{cohérent} relativement à un
  système~$\Sigma$, ou, comme nous le dirons, $\Sigma$-cohérent, si et
  seulement si $\Gamma \Proves/[\Sigma] \lfalse$.
\end{defn}

Par exemple, l'ensemble $\{ \Box(p \lif q), \Box p, \lnot\Box q \}$
est cohérent relativement à la logique propositionnelle, mais n'est
pas \Log{K}-cohérent. De même, l'ensemble
$\{ \Diamond p, \Box\Diamond p \lif q, \lnot q \}$ n'est pas
\Log{K5}-cohérent.

\begin{prop}\ollabel{prop:consistencyfacts}
  Soit $\Gamma$ un ensemble de !!{formula}s. Alors :
  \begin{enumerate}
  \item $\Gamma$ est $\Sigma$-cohérent si et seulement s'il existe
    une !!{formula}~$!A$ telle que $\Gamma \Proves/[\Sigma]
    !A$.
  \item \ollabel{prop:consistencyfacts-b}%
    $\Gamma \Proves[\Sigma] !A$ si et seulement si $\Gamma \cup \{
    \lnot!A \}$ n'est pas $\Sigma$-cohérent.
  \item \ollabel{prop:consistencyfacts-c}%
    Si $\Gamma$ est $\Sigma$-cohérent, alors, pour toute !!{formula}
    $!A$, au moins l'un des deux ensembles $\Gamma \cup \{ !A \}$
    et $\Gamma \cup \{ \lnot!A \}$ est $\Sigma$-cohérent.
  \end{enumerate}
\end{prop}

\begin{proof}
  Ces faits découlent aisément de la logique propositionnelle
  classique. Démontrons le point \olref{prop:consistencyfacts-c} par
  contraposition. Supposons que ni $\Gamma \cup \{ !A \}$ ni
  $\Gamma \cup \{ \lnot!A \}$ ne soient $\Sigma$-cohérents. Par la
  définition de la cohérence, on a alors à la fois
  $\Gamma, !A \Proves[\Sigma] \lfalse$ et
  $\Gamma, \lnot !A \Proves[\Sigma] \lfalse$. Le théorème de
  déduction donne $\Gamma \Proves[\Sigma] !A \to \lfalse$ et
  $\Gamma \Proves[\Sigma] \lnot!A \lif \lfalse$. Or
  $(!A \lif \lfalse) \lif ((\lnot!A \lif \lfalse) \lif \lfalse)$ est une
  instance de tautologie ; donc, d'après
  la \olref[prp]{prop:derivabilityfacts}\olref[prp]{prop:derivabilityfacts-ruleT},
  $\Gamma \Proves[\Sigma] \lfalse$.
\end{proof}
\end{document}
